\chapter{附录：构建、源码与模拟答案}

\section{编译 PDF}

\begin{verbatim}
cd /mnt/e/recoveryProjects/docs/engineering/cdc-two-gen-beginner
make pdf
# 输出: build/main.pdf
\end{verbatim}

需要 WSL 下 \texttt{texlive-xetex}、\texttt{texlive-lang-chinese}、\texttt{pgfplots}。

\section{源码文件（按需跳转）}

\begin{table}[htbp]
\centering
\small
\begin{tabular}{@{}ll@{}}
\toprule
\textbf{主题} & \textbf{路径} \\
\midrule
第一代 FastCDC & \filepath{engine\_cpp/src/chunk/fast\_cdc*.cc} \\
第二代 Build2FMasks & \filepath{engine\_cpp/src/chunk/gt\_cdc\_gear\_scan.cc} \\
第二代流式 & \filepath{engine\_cpp/src/chunk/gt\_cdc\_streaming.cc} \\
单元测试 & \filepath{engine\_cpp/tests/chunk/gt\_cdc\_2f\_test.cc} \\
玩具模拟数据 & \filepath{docs/engineering/cdc-two-gen-beginner/chapters/sim-toy-data.tex} \\
\bottomrule
\end{tabular}
\end{table}

\section{伪代码对照（两代并列）}

\begin{codebox}[第一代 vs 第二代 — 切点判断一行对比]
\begin{lstlisting}
// 第一代
if ((h & mask) == 0) cut;

// 第二代
if ((h_gear & mask_lo) == 0 &&
    ((h_norm >> norm_shift) & mask_hi) == 0) cut;
\end{lstlisting}
\end{codebox}

\section{模拟习题参考答案}

\begin{simbox}[Chunk 3（第一代，pos=7）]
暖机字节 \texttt{0x80, 0x90}：$h=20$；$(20 \ll 1)+11-20=31$。\\
$i=9$：$31\&3=3$ 不切；$h=(31 \ll 1)+6-20=48$。\\
$i=10$：$48\&3=0$ → 切点 10，Chunk 3 = 字节 7..9。
\end{simbox}

\begin{simbox}[两代切点对照（同一 12 字节文件）]
\begin{tabular}{@{}lll@{}}
\toprule
\textbf{Chunk} & \textbf{第一代切点} & \textbf{第二代切点} \\
\midrule
1 & 3 & 4 \\
2 & 7 & 11 附近（自练） \\
3 & 10 & （自练） \\
\bottomrule
\end{tabular}
\end{simbox}

\vspace{2em}
\noindent\textit{— 完 —}
